\documentclass[sigconf]{support/acmart}

% DepthLane paper scaffold.
% The original ACM sample files in this directory are intentionally unchanged.
% TODO before submission: replace venue metadata, author block, CCS terms, and all
% red TODO markers; remove forced page breaks after the content stabilizes.

\usepackage{booktabs}
\usepackage{multirow}

\setcopyright{none}
\settopmatter{printacmref=false, printfolios=true}
\renewcommand\footnotetextcopyrightpermission[1]{}

\acmConference[ACT Workshop '26]{ACT Workshop 2026}{2026}{TBD}
\acmBooktitle{Proceedings of ACT Workshop 2026}
\acmDOI{}
\acmISBN{}

\newcommand{\system}{\textsc{DepthLane}}
\newcommand{\todo}[1]{\textcolor{ACMRed}{[TODO: #1]}}
\newcommand{\figplaceholder}[2]{%
  \fbox{\parbox[c][#1][c]{0.96\linewidth}{\centering\small\bfseries #2}}}

\begin{document}

\title{DepthLane: Pipelined GPU--FPGA Split Inference for Real-Time Metric Depth Completion}

% TODO: replace with the venue's required anonymous or non-anonymous author block.
\author{Anonymous Author(s)}
\affiliation{%
  \institution{Anonymous Institution}
  \city{Anonymous}
  \country{Anonymous}}

\renewcommand{\shortauthors}{Anonymous Author(s)}

\begin{abstract}
Metric depth completion converts sparse or noisy range measurements into dense
metric geometry.  Running it directly on an embedded SoC keeps sensor data
local, avoids cloud round trips, and provides low-latency geometric information
for applications such as autonomous driving and mobile robotics.  However,
existing metric depth completion models are predominantly evaluated on
discrete GPUs, and few studies demonstrate high-frame-rate execution on
resource- and power-constrained embedded SoCs.  This leaves a practical gap
between model-level accuracy and real-time edge deployment.  We present
\system, a heterogeneous inference system that executes the RGB and RGB-D
encoders on an NVIDIA Jetson AGX Thor and the decoder on an Intel Stratix 10 MX
FPGA connected through four high-speed Ethernet lanes. Compared with
executing the complete model on Thor, \system{} achieves a 1.41$\times$
end-to-end speedup while maintaining the SOTA metric-depth
accuracy among other depth completion model. It also reduces total energy per output frame by \todo{XX\%} relative to the
all-on-Thor baseline.  The prototype further demonstrates that a complete model
need not reside on one processor: a compute-intensive module can be moved to an
FPGA behind a well-defined tensor interface, providing a reusable path for
accelerating compatible modules from other models.  
\end{abstract}

\keywords{heterogeneous computing, FPGA, edge inference, depth completion,
post-training quantization, split inference, frame pipelining}

\maketitle

% =============================== PAGE 1 ===============================
\section{Introduction}

Metric depth estimation is a key building block for spatial understanding,
robotics, measurement, and scene reconstruction.  Modern encoder--decoder
networks provide strong accuracy, but their computational and memory demands
make sustained real-time inference difficult on an embedded processor.  A
heterogeneous GPU--FPGA platform offers complementary resources: the GPU is
well suited to the irregular, high-level feature extraction of an encoder,
whereas the FPGA can implement the regular convolutional dataflow of a decoder.

This work studies AS-Depth as a concrete split-inference case.  The model
consumes RGB and sparse or noisy depth and produces a dense metric-depth map.
Its transformer-based encoder executes on an NVIDIA Jetson AGX Thor, while its
convolutional decoder executes on an Intel Stratix 10 MX FPGA.  The two devices
exchange intermediate features and depth outputs over four optical Ethernet
lanes.  The decoder is quantized to INT8 using PTQ so that it fits the FPGA
compute and memory budget without requiring retraining.

Simply moving the decoder to the FPGA does not guarantee higher throughput.
For a single frame, encoder execution, feature transfer, FPGA decoding, and
result return remain on the end-to-end critical path.  Across consecutive
frames, however, these stages use independent engines.  Once the features of
frame $i$ leave Thor, the GPU can encode frame $i+1$ while the FPGA decodes
frame $i$.  \system{} uses double buffering to turn this inter-frame
independence into a steady-state pipeline that hides much of the FPGA-side
latency.

This paper makes the following contributions:
\begin{itemize}
  \item We implement AS-Depth as a two-device pipeline: Thor executes the
  encoders, while the FPGA executes an INT8 PTQ decoder.  Overlapping the
  encoder of frame $i+1$ with the decoder of frame $i$ yields a
  1.41$\times$ end-to-end speedup over the all-on-Thor baseline while keeping
  relative metric-depth degradation below 5\%.
  \item We demonstrate the feasibility of module-level FPGA offload through a
  buffered tensor interface.  Because the controller depends on the module's
  inputs, outputs, and completion signal rather than AS-Depth semantics, it can
  be reused for compatible compute-intensive modules from other models.
  \item We evaluate power at the complete-system boundary, including Thor, the
  FPGA board, and the interconnect, and reduce energy per output frame by
  \todo{XX\%} relative to the all-on-Thor baseline.
\end{itemize}

\paragraph{Scope.}
We do not propose a new depth network or target a domain-specific application
stack.  AS-Depth is the case study used to evaluate a module-level
heterogeneous pipeline.  Generalization requires a valid model
partition, sufficient link bandwidth, and an FPGA-compatible offloaded stage;
Section~\ref{sec:limitations} states these limits.

% The full-width overview is allowed to float to the next available page.
\begin{figure*}[t]
  \centering
  \figplaceholder{1.42in}{%
    LEFT (about 62\%): system architecture.  RGB+sparse depth enters Thor;
    the AS-Depth encoder emits an INT8 feature stream over 4$\times$10GbE;
    Stratix 10 MX executes the decoder and returns uint16 metric depth.  Show
    ping/pong feature and output buffers.  Overlay frame $i$ on the FPGA path
    and frame $i+1$ on the GPU encoder to emphasize inter-frame overlap.\par
    RIGHT (about 38\%): aligned qualitative example with (a) RGB input,
    (b) sparse/noisy depth input, and (c) predicted dense metric depth.}
  \caption{\system{} overview.  Consecutive frames occupy different pipeline
  stages: Thor encodes frame $i+1$ while frame $i$ is transferred and decoded
  on the FPGA.}
  \Description{A planned system diagram showing an AS-Depth encoder on Thor,
  four optical Ethernet lanes, an INT8 decoder on a Stratix 10 MX FPGA, double
  buffers, and a qualitative metric-depth output.}
  \label{fig:overview}
\end{figure*}

\section{Background and Motivation}

\subsection{AS-Depth Split Point}

At the target input size of $518\times686$, the AS-Depth encoder produces a
$2048\times37\times49$ fused feature map before the convolutional decoder.
This boundary is attractive because it does not require transferring a full set
of U-Net skip tensors.  Static analysis of the current decoder gives
approximately 11.9 million parameters and 135 GMAC per frame (one
multiply--accumulate counted as one MAC).  The decoder expands the patch-grid
representation through multiple convolution and upsampling stages before
producing metric depth.

The split also exposes an important systems constraint.  The boundary tensor is
approximately 3.54 MiB in INT8.  Serialization alone is about 29.7 ms at
1 Gb/s, 2.97 ms at 10 Gb/s, and 0.74 ms at an ideal aggregate 40 Gb/s.  Thus,
the current 1 Gb/s software result is not sufficient for a low-latency system.
Before the accelerator evaluation, we will distinguish physical link rate from
host packet-generation limits and report effective payload throughput per lane
and in aggregate.  This communication-aware framing follows the central lesson
of earlier split-inference systems \cite{matsubara2019split}, while our focus is
the pipeline opportunity created by a local GPU--FPGA system.

\subsection{Inter-Frame Pipeline Opportunity}

A serialized split implementation processes $E_i\rightarrow X_i\rightarrow
F_i\rightarrow R_i$ before starting frame $i+1$, where $E$ is GPU encoding,
$X$ and $R$ are feature and result transport, and $F$ is FPGA decoding.  This
schedule leaves one device idle while the other is active.  With independent
buffers, $E_{i+1}$ can overlap $X_i$, $F_i$, and $R_i$.  The steady-state frame
interval is therefore governed by the slower pipeline side rather than the sum
of all stage times, although the latency of an individual frame still includes
every stage.  Prior FPGA systems similarly use layer- or task-level pipelining
to improve throughput \cite{razavi2022patterns,symons2022stream}; our focus is
frame-level overlap across an existing GPU--FPGA model boundary.

\subsection{Reusable Module Offload}

The runtime treats the offloaded decoder as a module with an input tensor
contract, an output tensor contract, and a completion event.  It does not depend
on depth-specific scheduling semantics.  The same transport and double-buffer
structure can therefore support another model partition if the selected module
can be implemented efficiently on the FPGA and its boundary tensors fit the
available link and buffer capacity.  AS-Depth provides a demanding case study,
not a restriction of the system architecture.

\subsection{Design Goals}

\system{} has four goals: (G1) limit relative metric-depth degradation after
PTQ to at most 5\%; (G2) sustain at least 30 frames/s in steady state; (G3)
include packing, communication, FPGA decoding, and result return in all
end-to-end measurements; and (G4) quantify energy and FPGA resource cost for
the complete pipeline.  We treat a missed target as a measured limitation
rather than reporting accelerator-kernel time alone.

\section{Heterogeneous System Design}
\label{sec:design}

\begin{figure*}[t]
  \centering
  \figplaceholder{1.05in}{%
    Execution timeline with rows for Thor encoder, optical TX/RX, FPGA decoder,
    and completed output.  Serialized split: Enc$_i$, TX$_i$, Dec$_i$, and
    RX$_i$ finish before Enc$_{i+1}$.  Pipelined split: Enc$_{i+1}$ overlaps
    TX$_i$+Dec$_i$+RX$_i$, followed by Enc$_{i+2}$ overlapping the offload path
    of frame $i+1$.  Mark ping/pong buffers, pipeline fill/drain, single-frame
    latency, and steady-state frame interval.}
  \caption{Inter-frame pipeline schedule.  Thor encodes the next frame while
  the current frame traverses the optical link and FPGA decoder.  This overlap
  targets throughput rather than removing stages from single-frame latency.}
  \Description{A planned Gantt chart comparing serialized and pipelined
  execution of consecutive frames across Thor, the optical link, and the FPGA.}
  \label{fig:timeline}
\end{figure*}

\subsection{INT8 Post-Training Quantization}

The decoder is quantized without retraining.  We use signed INT8 weights with a
separate symmetric scale for each output channel, signed INT8 activations with a
static per-tensor scale, and INT32 accumulation.  A calibration set of
\todo{128 or 256} training frames determines activation ranges.  We will
compare simple min--max calibration with \todo{percentile or MSE clipping} and
select one method before evaluating the held-out validation set.  Integer-only
inference follows standard affine requantization practice
\cite{jacob2018quantization}.

Every supported operation will have a bit-exact host reference.  The reference
fixes rounding, saturation, accumulator width, bias scaling, and requantization
before RTL validation.  If a fully INT8 decoder violates the predeclared
5\% relative-accuracy bound, the only fallback evaluated will be an INT16 final
projection; we will not introduce QAT or a learned compression module in this
paper.

\subsection{Streaming FPGA Decoder}

The decoder uses DSP-parallel convolution engines, INT32 partial sums, and
line/tile buffers for activation reuse, following the general streaming
accelerator approach exemplified by FINN \cite{umuroglu2017finn}.  Weights are loaded before the timed
run and are excluded only from per-frame latency when the same exclusion is
applied to GPU baselines.  Consecutive layers are consumed in producer--consumer
order when buffer capacity permits.  The stride-2, kernel-2 transposed
convolution can be implemented as an equivalent pointwise projection followed
by pixel shuffle, avoiding overlap handling in the upsampling datapath.

This paper will report post-fit ALMs, DSPs, M20K bits/blocks, achieved clock,
and any external-memory traffic.  The existing optical design consumes about
7\% of ALMs and no DSP blocks, leaving substantial arithmetic capacity, but the
final compute design must close timing; ``Fitter successful'' is not treated as
sufficient evidence.

\subsection{Optical Feature Transport}

Thor and the Stratix 10 MX exchange raw Ethernet frames over four 10GbE lanes.
Jumbo payloads amortize packet overhead.  Each packet carries a frame identifier,
lane/stripe identifier, byte offset, payload length, and checksum/status field.
Fixed-size striping distributes consecutive feature blocks across lanes.  The
FPGA writes arriving blocks into the current ping or pong buffer and exposes a
frame-ready event after all expected blocks pass validation.  Returned depth is
encoded as a $518\times686$ uint16 map unless the final preprocessing pipeline
selects a different padded resolution.

The transport is intentionally loss-intolerant for the evaluation: a missing or
invalid packet marks the frame incomplete and removes it from useful goodput.
We report offered traffic, received payload throughput, packet/frame loss, and
CPU overhead.  Retransmission and safety certification are future work.

\subsection{End-to-End Frame Pipeline}

Two feature buffers and two output buffers permit frame $i+1$ to enter the GPU
while frame $i$ is decoded.  Timestamps are recorded at input release, encoder
completion, final packet transmission, FPGA completion, result arrival, and
application delivery.  This distinguishes steady-state period from individual
frame latency and prevents an overlap-only throughput gain from being reported
as a single-frame latency reduction.  After pipeline fill, the approximate
frame interval is
\begin{equation}
  T_{\mathrm{frame}} \approx \max(T_E, T_{\mathrm{offload}}),
  \label{eq:pipeline-period}
\end{equation}
where $T_{\mathrm{offload}}$ includes feature transport, FPGA decoding, and
result return, including any overlap internal to that path.  Sustaining 30
frames/s requires $T_{\mathrm{frame}}\leq 33.3$\,ms.

\section{Inter-Frame Pipeline Execution}
\label{sec:pipeline}

\subsection{Stage and Buffer Model}

Each frame passes through GPU encoding $E_i$, feature transport $X_i$, FPGA
decoding $F_i$, and result return $R_i$.  Dependencies are enforced only within
the same frame.  Across frames, $E_{i+1}$ may begin after a free feature buffer
is available, without waiting for $F_i$ or $R_i$ to complete.  Ping/pong feature
and output buffers prevent the producer from overwriting data still consumed by
the link or FPGA.

\subsection{Pipeline Controller}

The host controller maintains buffer ownership and frame identifiers.  It
submits $E_{i+1}$ while the offload path processes frame $i$, then associates
each returned depth map with its original input.  Backpressure stops a new
encoder launch when both feature buffers are occupied.  This bounded policy
avoids unbounded queues while preserving every frame during the offline
evaluation.  The controller is dedicated to consecutive frames of the same
model and makes no multi-model admission decision.

\subsection{Generality Beyond AS-Depth}

The pipeline controller only requires tensor dimensions, serialization rules,
buffer addresses, and completion events for the selected partition.  A module
from another model can replace the AS-Depth decoder when (1) its operators map
efficiently to the FPGA, (2) its weights and working set fit local or HBM
capacity, and (3) boundary-transfer time does not dominate its saved compute.
The achievable speedup remains model- and split-dependent, so we present this as
a reusable system mechanism rather than a universal acceleration claim.

\subsection{Baselines}

We compare four configurations.  \emph{Thor FP32} and \emph{Thor TensorRT
FP16} execute the complete model on Thor.  \emph{Serialized split} executes the
Thor encoder and INT8 FPGA decoder without overlap between consecutive frames.
\emph{Pipelined split} uses identical hardware, quantization, and transport but
overlaps $E_{i+1}$ with the offload path of frame $i$.  Comparing the last two
isolates the benefit of inter-frame pipelining from model partitioning and PTQ.

\section{Prototype and Methodology}

Tables~\ref{tab:thor-hardware} and~\ref{tab:fpga-hardware} summarize the two
hardware platforms.  The FPGA table reports device capacity; post-fit resources
consumed by our decoder are reported separately in Table~\ref{tab:primary}(b).

\begin{table*}[t]
  \caption{NVIDIA Jetson AGX Thor hardware configuration.  Peak throughput and
  port capability are vendor specifications \cite{nvidiathorintro,nvidiathorspec};
  average power is measured at the platform input while running the evaluated
  workload.  The prototype operates the QSFP28 link at 4$\times$10\,GbE.}
  \label{tab:thor-hardware}
  \centering
  \small
  \begin{tabular}{p{0.30\textwidth}p{0.30\textwidth}p{0.32\textwidth}}
    \toprule
    CPU & GPU & Peak AI throughput \\
    \midrule
    14-core Arm Neoverse V3AE, 2.6\,GHz &
    2,560-core Blackwell, 10 TPCs &
    1,035 INT8 TOPS; 2,070 sparse-FP4 TFLOPS \\
    \midrule
    Supported precision & Average power & Optical interface \\
    \midrule
    FP32/TF32/BF16/FP16/FP8/FP4; INT8; INT4 weights &
    \todo{measured W} (40--130\,W mode range) &
    1$\times$QSFP28 (4$\times$25\,GbE); used at 4$\times$10\,GbE \\
    \bottomrule
  \end{tabular}
\end{table*}

\begin{table*}[t]
  \caption{Intel Stratix 10 MX FPGA hardware configuration.  Values are total
  device capacities rather than post-fit utilization \cite{intels10mxproduct};
  Intel reports ALMs rather than LUTs.  Average power is measured at the board
  input, and the two QSFP28 cages are 100GbE-capable while this prototype uses
  four 10\,GbE lanes \cite{intels10mxboard,intels10mxdsp}.}
  \label{tab:fpga-hardware}
  \centering
  \small
  \begin{tabular}{p{0.30\textwidth}p{0.30\textwidth}p{0.32\textwidth}}
    \toprule
    Device & Logic resources & DSP resources and precision \\
    \midrule
    1SM21BHU2F53E2VGS1 (MX 2100) &
    2,073,000 LEs; 702,720 ALMs; 2,810,880 registers &
    3,960 DSPs; 2$\times$18$\times$19 or 1$\times$27$\times$27 fixed-point; hardened FP32 \\
    \midrule
    Memory resources & Average power & Optical interface \\
    \midrule
    6,847 M20Ks (134\,Mb); 11\,Mb MLAB; 94.5\,Mb eSRAM; 8\,GB HBM2 &
    \todo{measured W} &
    2$\times$QSFP28 (100GbE capable); used at 4$\times$10\,GbE \\
    \bottomrule
  \end{tabular}
\end{table*}

All input files are decoded before a timed run and replayed from memory.  This
removes storage variability without requiring a live sensor, while preserving
real model execution, packet transfer, buffering, and synchronization.  Each
sample supplies the aligned RGB and sparse/noisy-depth inputs expected by
AS-Depth.  Accuracy runs process every sample without overlap-dependent drops;
throughput runs replay the same fixed input sequence for every baseline.

Thor power is sampled from the platform's exposed GPU and system-input rails
\cite{nvidiathorpower}.  FPGA power is measured at the development-board input
and cross-checked with the board tool where possible \cite{intels10mx}.  The
headline energy metric includes both devices and optical transceivers.  We also
report active-minus-idle energy to explain static board overhead, but do not use
it as a substitute for total system energy.

\section{Evaluation Plan and Results}
\label{sec:evaluation}

The evaluation answers three questions.  \textbf{RQ1:} Does INT8 decoder PTQ
keep metric-depth degradation within 5\%?  \textbf{RQ2:} Does the complete
Thor--FPGA pipeline sustain at least 30 frames/s, and at what latency, energy,
and FPGA cost?  \textbf{RQ3:} How much offload-path time is hidden by
inter-frame overlap relative to a serialized split?

\begin{table*}[t]
  \caption{Primary accuracy, systems, and FPGA results.  All numerical cells
  are intentionally blank until reproduced on the final software/bitstream.}
  \label{tab:primary}
  \centering
  \small
  \textbf{(a) Depth accuracy and end-to-end performance on Hammer}\par\smallskip
  \resizebox{0.96\textwidth}{!}{%
  \begin{tabular}{lrrrrrrrr}
    \toprule
    Configuration & RMSE (mm)$\downarrow$ & MAE (mm)$\downarrow$ & iRMSE (1/km)$\downarrow$ & iMAE (1/km)$\downarrow$ & p50 (ms)$\downarrow$ & p99 (ms)$\downarrow$ & FPS$\uparrow$ & J/frame$\downarrow$ \\
    \midrule
    Thor full model FP32 & -- & -- & -- & -- & -- & -- & -- & -- \\
    Thor full model TensorRT FP16 & -- & -- & -- & -- & -- & -- & -- & -- \\
    Serialized Thor encoder + FPGA INT8 decoder & -- & -- & -- & -- & -- & -- & -- & -- \\
    Pipelined Thor encoder + FPGA INT8 decoder & -- & -- & -- & -- & -- & -- & -- & -- \\
    \bottomrule
  \end{tabular}}

  \vspace{0.6em}
  \textbf{(b) Post-fit FPGA implementation}\par\smallskip
  \resizebox{0.82\textwidth}{!}{%
  \begin{tabular}{lrrrrrr}
    \toprule
    Design & ALMs & Registers & DSPs & Memory bits / M20Ks & Fmax (MHz) & Decoder latency (ms) \\
    \midrule
    Four-lane optical baseline & 47,261 & 54,618 & 0 & 704,064 / 77 & -- & N/A \\
    Optical + INT8 decoder & -- & -- & -- & -- & -- & -- \\
    \bottomrule
  \end{tabular}}
\end{table*}

\subsection{RQ1: Accuracy Preservation}

We evaluate the complete held-out Hammer set using RMSE, MAE, iRMSE, and iMAE.
The principal comparison uses the same checkpoint and preprocessing with a
floating-point decoder and the INT8 PTQ decoder.  For each error metric $m$, we
report $(m_{\mathrm{INT8}}-m_{\mathrm{float}})/m_{\mathrm{float}}$ and require
the relative increase to remain at or below 5\%.  This threshold and the
calibration procedure are fixed before evaluating the held-out set.  One
qualitative row in Fig.~\ref{fig:overview} will show a representative sample and
a difficult failure case if space permits.

\subsection{RQ2: Real-Time Throughput and System Cost}

Table~\ref{tab:primary} reports steady-state FPS, end-to-end latency, energy,
and post-fit resource use.  Real-time operation is supported only if the
pipelined configuration sustains at least 30 complete outputs/s, corresponding
to a frame interval no greater than 33.3\,ms.  We separately report encoder,
feature-transfer, decoder, return, and synchronization times so the source of
any missed target is visible.  Energy per frame is the integral of complete
platform power over the run divided by the number of returned depth maps.

\begin{figure}[t]
  \centering
  \figplaceholder{1.35in}{%
    Three aligned plots for Thor FP32, Thor FP16, serialized split, and
    pipelined split: (a) steady-state FPS with a 30-FPS target line; (b) encoder,
    transfer, FPGA decoder, return, and exposed-versus-hidden time; (c) average
    watts and joules per frame.  Use identical input sequences and measured
    confidence intervals.}
  \caption{Planned throughput, overlap, and energy comparison.  Serialized and
  pipelined split execution use the same model, INT8 decoder, hardware, and
  transport configuration, isolating the benefit of inter-frame overlap.}
  \Description{A placeholder for throughput, stage-latency, hidden-time, and
  energy plots comparing full-model Thor execution with serialized and
  pipelined GPU--FPGA split inference.}
  \label{fig:pipeline-results}
\end{figure}

\subsection{RQ3: Pipeline Overlap Effectiveness}

Figure~\ref{fig:pipeline-results} compares serialized and pipelined split
execution over the same ordered frame sequence.  We report pipeline fill and
drain separately from steady state, together with the fraction of offload-path
time hidden behind $E_{i+1}$.  The pipeline claim is supported only if overlap
improves completed-frame throughput without changing output tensors or dropping
evaluation frames.  We do not claim that a 30-FPS throughput implies a
single-frame latency below 33.3\,ms; these are reported as distinct quantities.

\section{Discussion and Limitations}
\label{sec:limitations}

The prototype evaluates prerecorded inputs rather than live sensor I/O.  The
observed accuracy preservation is specific to the evaluated checkpoint,
calibration set, and decoder; PTQ does not guarantee the same bound for another
model.  Insufficient effective link throughput can dominate the split even when
the FPGA decoder is fast.  The 30-FPS target concerns steady-state throughput
after pipeline fill, not the latency of an isolated frame.  Generalization to
another model is also not automatic: the selected module must have supported
operators, adequate parallelism, a manageable boundary tensor, and a working
set that fits the FPGA memory hierarchy.  Activation compression, dynamic
repartitioning, and fault recovery are outside the present scope.

\section{Conclusion}

\system{} is a GPU--FPGA split-inference system for real-time metric depth
estimation.  Thor executes the AS-Depth encoder, while a PTQ INT8 decoder runs
on the FPGA.  Double buffering overlaps encoder execution for the next frame
with transport and FPGA decoding of the current frame, hiding offload-path time
in steady state.  Relative to executing the complete model on Thor, this design
achieves a 1.41$\times$ end-to-end speedup while keeping relative metric-depth
degradation below 5\%, and reduces energy per output frame by \todo{XX\%}.
Beyond AS-Depth, the buffered module interface provides a path to offload
compatible stages from other models, subject to operator, memory, and
communication constraints.

% Flush the deferred full-width results table before starting the final
% reference page.  Earlier sections deliberately contain no forced breaks.
\clearpage
\bibliographystyle{support/ACM-Reference-Format}
\bibliography{references/depthlane}

\end{document}
